\subsection{Robust candidate and method selection}

The strongest immediate extension is a sequence-aware candidate selector.
Current oracle and method-portfolio experiments show that plausible boundaries
are often present in the candidate pool, but the system does not reliably pick
the right subset for each video. A better selector should combine local
semantic contrast, cross-model agreement, pause/shot/OCR evidence, position
priors, and global duration constraints while optimising directly against
$P_k$, WindowDiff, and tolerance-F1.

\subsection{Independent validation of chapter references}

The most important dataset extension is not simply adding more videos; it is
checking whether creator-provided YouTube chapters agree with independent
human judgements. A focused validation study should select 8--10 lectures and
ask two annotators to mark chapter-level boundaries without seeing the YouTube
timestamps. The study should report human--human agreement, human--YouTube
agreement, and system--human agreement under F1@2, F1@5, and time-tolerant
metrics. This would quantify how reliable the current creator chapters actually
are as reference boundaries.

\subsection{Expanded annotator pool and domain diversity}

The current inter-annotator study involves annotators from a shared academic
background. Future work should recruit annotators from different disciplines
and experience levels to test whether subtopic agreement generalises beyond
the current pool. Annotating 8--10 additional videos with three or more
independent annotators would give a Krippendorff's $\alpha$ estimate alongside
the existing Cohen's $\kappa$, and would allow pairwise agreement patterns to
be analysed rather than collapsed to a single mean figure.

\subsection{External and Math-focused validation}

The current selector fails most clearly on Mathematics, where the dataset has
only four videos. A targeted extension should add 8--10 Math lectures with
creator chapters and re-evaluate the final methods without retuning. If Math
performance improves, the likely cause was training scarcity; if it stays weak,
the problem probably runs deeper, involving mathematical notation, stable
vocabulary, or gradual conceptual transitions. A separate 5--10 video external
test set from new channels should also be used to measure generalisation beyond
\dataset{}.

\subsection{Reliability-aware prosody and visual integration}

Prosodic features (pause duration, pitch reset) and visual streams (shot
boundaries, slide OCR) are available in the pipeline, but current ablations
show that acoustic and linguistic sub-sentence signals can hurt
$P_k$/WindowDiff when their granularity does not match chapter-level editorial
decisions. The exception is CLIP visual embeddings: CLIP alone achieves
$P_k=0.3958$ and CLIP plus text fusion reaches $P_k=0.3740$, approaching the
cross-model text method ($P_k=0.3715$). Visual scene changes in lecture slides
appear to carry chapter-level granularity. Future work should prioritise CLIP
visual integration within the cross-model scoring framework, and apply
quality-gating only to the noisier acoustic and OCR modalities.

\subsection{Streaming and online inference}

The current pipeline processes a complete video offline.
A streaming variant would segment lecture content as it is being recorded,
enabling real-time chaptering for live classes.
This requires adapting the sliding-window fusion to operate on a causal context
window and replacing the global threshold computation with an exponential moving
average of past gap scores.
The target latency budget for live chaptering is under 30~seconds of delay.

\subsection{Cross-lingual extension}

\lecseg{} currently targets English lectures.
Extending to other languages requires: (1)~replacing \texttt{all-mpnet-base-v2}
with a multilingual text encoder such as LaBSE or \texttt{multilingual-e5-large};
and (2)~using Whisper's multilingual mode for transcription.
Suitable evaluation data could come from non-English MOOCs (e.g.,
Coursera's French and Spanish tracks) or from university recordings in South
and South-East Asian languages.

\subsection{End-to-end fine-tuning}

The current architecture uses frozen pre-trained encoders.
Fine-tuning the text encoder jointly with the boundary predictor on a larger
annotated corpus (e.g., constructed by scaling the \dataset{} annotation
procedure to 200+ videos) could yield further gains, at the cost of reduced
generalisability.
A natural training objective is a soft version of the tolerance-$F_1$ metric,
differentiable via a sigmoid boundary-probability output.

\subsection{Larger and domain-tuned local LLMs}

Replacing Llama-3.1-8B with a 70B parameter model or a domain-specific
fine-tune (e.g., fine-tuned on lecture transcripts) could improve title
quality and boundary-shift accuracy, particularly for highly technical content.
Quantised models (GGUF 4-bit) make 70B inference feasible on a single 24\,GB
GPU with acceptable latency (around 5~seconds per boundary refinement call).

\subsection{Multi-speaker lectures and panels}

Panel discussions and multi-instructor courses introduce speaker turns as an
additional boundary signal.
Speaker diarisation (e.g., using pyannote.audio) could provide a speaker-change
feature that complements text embeddings, particularly for sections where the
topic does not change but the speaker does.

\subsection{User study on learning outcomes}

All metrics in this thesis are information-retrieval proxies for segmentation
quality.
A controlled user study would directly measure whether learners locate
target concepts faster, and whether their retention improves, when using
\lecseg{}-generated chapters compared to flat video scrubbing or manually
authored chapters.
Such a study could also reveal which granularity (chapter vs.\ subtopic)
is most useful for different learning tasks (review vs.\ first-watch).

\subsection{Public benchmark and model release}

A complete public release should include the manifest, video identifiers,
creator chapter references, reviewed subtopic labels, preprocessing scripts,
evaluation scripts, environment lock files, result-regeneration commands, and
dataset/model cards. The current repository already contains the major
components; the remaining work is release hygiene, privacy review, and a
single documented reproduction command for the final thesis tables.
